\documentclass[12pt]{article}
\usepackage{amsmath,amssymb,geometry,graphicx}
\usepackage[hidelinks]{hyperref}
\geometry{margin=1in}
\setlength{\parskip}{0.9ex}
\setlength{\parindent}{0pt}

\title{A Hypothesis for a Sub-Planckian Timescale Based on the Holographic Entropy of the Cosmos}
\author{Naser Ahani \\ Department of Physics, Amirkabir University of Technology \\ (Tehran Polytechnic), Tehran, Iran \\ \texttt{naser.ahani@aut.ac.ir}}
\date{November 2025}

\begin{document}
\maketitle

\begin{abstract}
This paper proposes that physical reality may possess a deterministic substructure operating on a timescale smaller than the Planck time. By relating this hypothetical sub-Planckian scale to the total holographic information content of the observable universe, a characteristic interval---the \textit{Planck-pulse timescale} ($T_{pp}$)---is derived. Using the holographic principle and the cosmological constant $\Lambda$ as the defining boundary of the universe, we obtain a universal expression for the total informational degrees of freedom, $N_Z \approx 3.33 \times 10^{122}$, yielding $T_{pp} \approx 1.62 \times 10^{-166}\,\mathrm{s}$. The result is discussed in the context of emergent time and information-theoretic interpretations of spacetime, and limitations and possible routes toward a statistical reconstruction of quantum mechanics are outlined.
\end{abstract}

\section{Introduction}
The incompatibility between General Relativity (GR) and Quantum Mechanics (QM) remains a central problem in fundamental physics. GR describes spacetime as a smooth dynamical manifold while QM assigns probabilistic amplitudes to measurement outcomes. Reconciling the two frameworks has motivated multiple approaches (string theory, loop quantum gravity, causal set theory, holographic models), each emphasizing different aspects of the putative microstructure of spacetime \cite{tHooft1993,Susskind1995,Rovelli2004}.

A recurring idea in recent decades is that spacetime and gravity may be emergent, arising from more primitive information-theoretic or thermodynamic degrees of freedom \cite{Jacobson1995,Padmanabhan2010,Verlinde2011}. Within that perspective, notions such as horizon entropy and holography provide quantitative footholds for counting microstates. This paper explores a concrete, minimal hypothesis along these lines: that there exists a deterministic causal update timescale below the Planck time, whose effective coarse-graining produces the familiar quantum statistical behavior.

Two constraints guide the construction: (i) keep the argument minimal and anchored to established relations (Bekenstein–Hawking entropy, de Sitter horizon determined by $\Lambda$); (ii) preserve dimensional consistency and explicit numerical estimates using standard physical constants. The main aim is to derive an expression for the number of informational degrees of freedom ($N_Z$) of the observable universe and, from it, a candidate sub-Planckian timescale $T_{pp}$.

\section{Theoretical Framework}
\subsection{Holography and horizon entropy}
Bekenstein and Hawking showed that the entropy of a black hole horizon is proportional to its area measured in Planck units:
\begin{equation}\label{eq:SBH}
S_{BH} = \frac{k_B A}{4 l_p^2}, \qquad l_p^2 \equiv \frac{\hbar G}{c^3}.
\end{equation}
The holographic principle extends this idea, suggesting that the maximal information content in a spatial region is bounded by the area of its boundary \cite{tHooft1993,Susskind1995}.

For a universe whose large-scale geometry is effectively de Sitter due to a positive cosmological constant $\Lambda$, the cosmological event horizon radius is
\begin{equation}\label{eq:RH}
R_H = \sqrt{\frac{3}{\Lambda}},
\end{equation}
with corresponding horizon area
\begin{equation}\label{eq:AH}
A_H = 4\pi R_H^2 = \frac{12\pi}{\Lambda}.
\end{equation}

\subsection{Definition of $N_Z$ and the Planck-pulse timescale}
We define $N_Z$ as the total number of fundamental information units (``bits'') that can be stored on the cosmic holographic screen. Taking one information unit to correspond to an area $4 l_p^2$ (consistent with the Bekenstein–Hawking factor 1/4), we obtain:
\begin{equation}\label{eq:NZ}
N_Z \equiv \frac{A_H}{4 l_p^2} = \frac{3\pi c^3}{\Lambda \hbar G}.
\end{equation}
This expression is dimensionless and time-independent so long as $\Lambda$ is treated as a constant.

We then postulate that the fundamental deterministic timescale $T_{pp}$ (Planck-pulse timescale) is related to the Planck time $T_p$ by
\begin{equation}\label{eq:Tpp_def}
T_{pp} = \frac{T_p}{N_Z},
\end{equation}
where
\begin{equation}\label{eq:Tp}
T_p = \sqrt{\frac{\hbar G}{c^5}}.
\end{equation}
Eq.~\eqref{eq:Tpp_def} is introduced here as the axiom of the present framework: $T_p$ is interpreted as an emergent statistical unit while $T_{pp}$ is the putative ontological tick of the underlying deterministic substrate.

\section{Results and Discussion}
\subsection{Numerical evaluation}
We substitute standard CODATA-style values (rounded for clarity) into Eq.~\eqref{eq:NZ}:
\[
c \approx 3.00\times10^{8}\,\mathrm{m\,s^{-1}},\quad
\hbar \approx 1.054\times10^{-34}\,\mathrm{J\,s},\quad
G \approx 6.674\times10^{-11}\,\mathrm{m^3\,kg^{-1}\,s^{-2}},
\]
and adopt the cosmological constant estimate inferred from Planck data:
\[
\Lambda \approx 1.088\times10^{-52}\,\mathrm{m^{-2}}.
\]
Thus
\begin{equation}\label{eq:NZ_numeric}
\begin{split}
N_Z &= \frac{3\pi (3.00\times10^8\,\mathrm{m\,s^{-1}})^3}
{(1.088\times10^{-52}\,\mathrm{m^{-2}})\,(1.054\times10^{-34}\,\mathrm{J\,s})\,(6.674\times10^{-11}\,\mathrm{m^3\,kg^{-1}\,s^{-2}})}\\
&\approx 3.33\times10^{122}.
\end{split}
\end{equation}
Using Eq.~\eqref{eq:Tp} we estimate the Planck time
\begin{equation}\label{eq:Tp_numeric}
T_p \approx \sqrt{\frac{1.054\times10^{-34}\times 6.674\times10^{-11}}{(3.00\times10^8)^5}} \approx 5.38\times10^{-44}\,\mathrm{s},
\end{equation}
and therefore
\begin{equation}\label{eq:Tpp_numeric}
T_{pp} = \frac{T_p}{N_Z} \approx \frac{5.38\times10^{-44}\,\mathrm{s}}{3.33\times10^{122}} \approx 1.62\times10^{-166}\,\mathrm{s}.
\end{equation}

\noindent These numerical evaluations are consistent with the order-of-magnitude estimates given in prior literature on cosmic entropy \cite{Egan2010} and illustrate the dramatic separation between the Planck timescale and the hypothesised ontological tick.

\subsection{Interpretation and conceptual consequences}
The interpretation adopted here is deliberately minimal: we do not claim to have derived a microscopic Hamiltonian or a unique underlying dynamics. Rather, the numerical result suggests the following conceptual picture:
\begin{itemize}
  \item The macroscopic Planck time $T_p$ may be a coarse-grained timescale resulting from averaging over $N_Z$ microscopic updates.
  \item Quantum indeterminacy could be an emergent, statistical phenomenon arising from an observer's limited access to sub-Planckian state variables.
  \item Time itself may be best understood as a measure of information change (updates per unit), with $T_{pp}$ the elementary informational update.
\end{itemize}
This picture resonates with thermodynamic and information-theoretic approaches to gravity \cite{Jacobson1995,Padmanabhan2010,Verlinde2011} but differs in proposing a concrete numerical separation between micro- and mesoscopic temporal resolution.

\subsection{Assumptions, limitations and sensitivity analysis}
Several assumptions and idealisations underpin the argument:
\begin{enumerate}
  \item \textbf{Horizon choice:} We used the de Sitter horizon set by $\Lambda$. Alternative choices (e.g., the particle horizon or the Hubble radius $c/H_0$) change numerical factors but not the qualitative scaling $N_Z\propto 1/\Lambda$.
  \item \textbf{Planck-area bit:} We associated one information unit with area $4l_p^2$ consistent with Bekenstein–Hawking prefactors. Different choices amount to $\mathcal{O}(1)$ rescalings of $N_Z$.
  \item \textbf{Constancy of $\Lambda$:} Treating $\Lambda$ as constant yields time-independent $N_Z$. If $\Lambda$ evolves (e.g., dynamical dark energy), $N_Z$ becomes time-dependent and the interpretation of $T_{pp}$ must be modified.
\end{enumerate}

A simple sensitivity estimate follows immediately from Eq.~\eqref{eq:NZ}: since $N_Z\propto 1/\Lambda$,
\begin{equation}
\frac{\delta N_Z}{N_Z} = -\frac{\delta\Lambda}{\Lambda}.
\end{equation}
Therefore, the relative uncertainty in $N_Z$ is the same order as the fractional uncertainty in the measured $\Lambda$. Given current cosmological constraints (Planck), this does not qualitatively affect the enormous scale separation reported above.

\subsection{Observational and theoretical prospects}
Direct experimental probes of processes at $T_{pp}$ are obviously impossible with present or foreseeable technology. Nevertheless, two indirect avenues may be considered:
\begin{itemize}
  \item \textbf{Theoretical reconstruction:} Show that ensembles of deterministic micro-updates with timescale $T_{pp}$ reproduce the statistical structure of quantum amplitudes at accessible scales. Constructing a toy model (cellular-automaton-like) or stochastic coarse-graining procedure is a plausible route.
  \item \textbf{Consistency conditions:} Demand that any underlying micro-dynamics respecting causality and locality at emergent scales reproduces known semiclassical relations (energy–entropy balances, black hole thermodynamics). This constrains candidate models even without direct measurement.
\end{itemize}

\section{Conclusions and Outlook}
We presented a simple derivation linking the holographic information content of the observable universe to a putative deterministic sub-Planckian timescale $T_{pp}$. The key results are compactly captured by Eqs.~\eqref{eq:NZ_numeric} and \eqref{eq:Tpp_numeric}. While speculative, this viewpoint offers a quantitatively explicit target for models that aim to reconstruct quantum statistical behavior from deterministic microphysics.

Future work should aim to:
\begin{itemize}
  \item Build explicit toy dynamical models that operate at $T_{pp}$ and investigate their coarse-grained quantum-like statistics;
  \item Study implications of a time-dependent $\Lambda$ for $N_Z$ and whether cosmological evolution leaves an imprint on emergent temporal statistics;
  \item Explore information-theoretic and causal-set approaches that naturally accommodate a large finite $N_Z$.
\end{itemize}

\vspace{1ex}
\noindent \textbf{Acknowledgement.} (none in this version)

\bibliographystyle{unsrt}
\begin{thebibliography}{99}
\bibitem{tHooft1993} 't Hooft, G. (1993). Dimensional Reduction in Quantum Gravity. \textit{arXiv:gr-qc/9310026}.
\bibitem{Susskind1995} Susskind, L. (1995). The World as a Hologram. \textit{Journal of Mathematical Physics}, 36(11), 6377--6396.
\bibitem{Penrose2004} Penrose, R. (2004). \textit{The Road to Reality: A Complete Guide to the Laws of the Universe}. Knopf.
\bibitem{Bekenstein1973} Bekenstein, J. D. (1973). Black Holes and Entropy. \textit{Phys. Rev. D}, 7(8), 2333--2346.
\bibitem{Jacobson1995} Jacobson, T. (1995). Thermodynamics of Spacetime: The Einstein Equation of State. \textit{Phys. Rev. Lett.}, 75, 1260--1263.
\bibitem{Padmanabhan2010} Padmanabhan, T. (2010). Thermodynamical Aspects of Gravity: New Insights. \textit{Reports on Progress in Physics}, 73, 046901.
\bibitem{Verlinde2011} Verlinde, E. (2011). On the Origin of Gravity and the Laws of Newton. \textit{JHEP}, 2011(4), 29.
\bibitem{Rovelli2004} Rovelli, C. (2004). \textit{Quantum Gravity}. Cambridge University Press.
\bibitem{Planck2020} Planck Collaboration (2020). Planck 2018 results. VI. Cosmological parameters. \textit{A\&A}, 641, A6.
\bibitem{Egan2010} Egan, C. A., \& Lineweaver, C. H. (2010). A Larger Estimate of the Entropy of the Universe. \textit{Astrophys. J.}, 710(2), 1825--1834.
\bibitem{Henson2009} Henson, J. (2009). The causal set approach to quantum gravity. In \textit{Approaches to Quantum Gravity}. Cambridge University Press.
\end{thebibliography}

\end{document}

